\section{Problem Definition} 
\label{sec:problem}

\subsection{VRPTW-S formulation}
The Vehicle Routing Problem with Time Windows and Synchronization constraints (VRPTW-S) is defined on a directed graph $G=(V,E)$, where $V={0,1,\dots,n}$ denotes the set of vertices and $E$ the set of edges. Vertex $0$ represents the depot, while each customer $i \in V\setminus{0}$ requires a service of duration $s_i$ to be performed within a time window $[e_i,l_i]$.

A fleet of vehicles is available to serve all customers. Traveling from customer $i$ to customer $j$ incurs a travel time $t_{ij}$ and an associated routing cost $c_{ij}$. Each customer must be visited exactly once, and every route must start and end at the depot.

In addition to classical routing and scheduling constraints, a subset of customers requires synchronized service. Let $P^{sync}$ denote the set of synchronized visit pairs. For each pair $(i,j)\in P^{sync}$, the service start times must satisfy
\[ T_i = T_j \]
where $T_i$ and $T_j$ denote the service start times at customers $i$ and $j$, respectively.

The objective is to determine a set of feasible routes minimizing the total routing cost while satisfying vehicle capacity, route duration, time-window, and synchronization constraints. Synchronization requirements create strong temporal dependencies between routes and significantly increase the complexity of feasibility management during optimization.

\subsection{Benchmark instances}
The experimental study relies on two complementary benchmark families.

The first consists of the synchronized Solomon-derived benchmark instances commonly used in the VRPTW-S literature. These instances originate from the benchmark generation procedures introduced by \citet{bredstrom2008combined} and subsequently extended by several authors. They provide the standard basis for evaluating and comparing synchronization-aware optimization methods.

The second benchmark family was developed specifically for this work to evaluate scalability beyond the instance sizes commonly considered in the literature. Starting from classical Solomon instances, additional synchronized visits were generated following the synchronization generation principles used in existing benchmark sets. The resulting dataset includes instances containing up to 400 customers and substantially higher synchronization densities than those found in traditional benchmarks.

While the literature benchmarks are used to validate the proposed GVNS+ implementation against previously reported results, all learning-guided experiments are conducted on the newly generated benchmark family. This experimental protocol allows the adaptive mechanisms to be evaluated independently from existing benchmark-specific tuning effects while providing a more challenging assessment of scalability and robustness.
